\documentclass{article}
\usepackage{amsmath,amssymb,amsthm}
\usepackage{algorithm,algorithmic}
\usepackage{enumitem}
\usepackage{hyperref}
\newtheorem{theorem}{Theorem}
\newtheorem{lemma}{Lemma}
\newtheorem{assumption}{Assumption}
\newtheorem{remark}{Remark}
\newcommand{\prox}{\operatorname{prox}}
\newcommand{\E}{\mathbb{E}}
\newcommand{\R}{\mathbb{R}}

\begin{document}

\section*{Algorithm Name}
\textbf{Accelerated Proximal Gradient with Momentum (APG-M)}  

The method solves the composite convex problem  

\[
\min_{x\in\R^{d}} \; \Phi(x)\;\stackrel{\text{def}}{=}\; f(x)+g(x),
\]

where $f$ is convex and $L$‑smooth, $g$ is convex (possibly non‑smooth) and its proximal operator is easy to compute.

--------------------------------------------------------------------

\section*{Mathematical Formulation}
Let $\eta>0$ be a stepsize satisfying $\eta\le 1/L$.  
Define the sequence $\{t_k\}_{k\ge 0}$ by  

\[
t_0\;=\;1,\qquad 
t_{k+1}\;=\;\frac{1+\sqrt{1+4t_k^{2}}}{2}, \qquad 
\beta_k\;=\;\frac{t_k-1}{t_{k+1}} .
\tag{1}
\]

The iterates are generated as  

\[
\boxed{
\begin{aligned}
y_k &= x_k + \beta_k\,(x_k-x_{k-1}), \qquad (k\ge 0,\;x_{-1}=x_0)\\[4pt]
x_{k+1} &= \prox_{\eta g}\bigl(y_k-\eta\nabla f(y_k)\bigr).
\end{aligned}}
\tag{2}
\]

When $g\equiv0$, (2) reduces to Nesterov’s accelerated gradient method; when $\beta_k\equiv0$, we obtain the standard proximal gradient method.

--------------------------------------------------------------------

\section*{Pseudocode}
\begin{algorithm}[H]
\caption{Accelerated Proximal Gradient with Momentum (APG-M)}
\begin{algorithmic}[1]
\Require Convex $L$‑smooth $f$, convex $g$, stepsize $\eta\in(0,1/L]$, initial point $x_0\in\R^d$, set $x_{-1}=x_0$, $t_0=1$.
\For{$k=0,1,2,\dots$}
    \State $t_{k+1}\gets \frac{1+\sqrt{1+4t_k^{2}}}{2}$ \hfill\Comment{update Nesterov coefficient}
    \State $\beta_k \gets \frac{t_k-1}{t_{k+1}}$
    \State $y_k \gets x_k + \beta_k (x_k-x_{k-1})$
    \State $x_{k+1} \gets \prox_{\eta g}\bigl(y_k-\eta\nabla f(y_k)\bigr)$
\EndFor
\end{algorithmic}
\end{algorithm}

Termination can be based on a maximum number of iterations $K$ or on the criterion $\|x_{k+1}-x_k\|\le\epsilon$.

--------------------------------------------------------------------

\section*{Dry Run Example}
Consider the scalar problem  

\[
f(w)=(w-3)^2,\qquad g(w)\equiv0,
\]
so that $\nabla f(w)=2(w-3)$, $L=2$, and $\prox_{\eta g}= \mathrm{id}$.  
Take $\eta=0.1$, $x_{-1}=x_0=0$.

\begin{center}
\begin{tabular}{c|c|c|c|c}
$k$ & $t_k$ & $\beta_k$ & $y_k$ & $x_{k+1}$ \\ \hline
0 & $1$ & $0$ & $0$ & $0-\eta\cdot2(0-3)=0.6$ \\[2pt]
1 & $\frac{1+\sqrt{5}}{2}\approx1.618$ & $\frac{0}{1.618}=0$ & $0.6$ & $0.6-\eta\cdot2(0.6-3)=0.6-0.1\cdot4.8=0.12$ \\[2pt]
2 & $\frac{1+\sqrt{1+4(1.618)^2}}{2}\approx2.193$ & $\frac{1.618-1}{2.193}\approx0.282$ & 
$0.12+0.282(0.12-0.6)= -0.0456$ &
$-0.0456-\eta\cdot2(-0.0456-3)= -0.0456-0.1\cdot(-6.0912)=0.5645$
\end{tabular}
\end{center}

Objective values after each iteration:  

\[
\Phi(x_1)= (0.6-3)^2=5.76,\quad
\Phi(x_2)= (0.12-3)^2=8.2944,\quad
\Phi(x_3)= (0.5645-3)^2\approx5.968.
\]

The sequence oscillates but quickly approaches the minimiser $w^\star=3$.

--------------------------------------------------------------------

\section*{Assumptions}
\begin{assumption}[Smoothness of $f$]\label{ass:smooth}
$f:\R^{d}\rightarrow\R$ is convex and differentiable with $L$‑Lipschitz gradient, i.e.  

\[
\|\nabla f(x)-\nabla f(y)\| \le L\|x-y\|,\qquad\forall x,y\in\R^{d}.
\tag{A1}
\]
\end{assumption}

\begin{assumption}[Convexity of $g$]\label{ass:convex}
$g:\R^{d}\rightarrow\R\cup\{+\infty\}$ is proper, closed and convex. Its proximal operator  

\[
\prox_{\eta g}(z)\;=\;\arg\min_{x}\Bigl\{g(x)+\frac{1}{2\eta}\|x-z\|^{2}\Bigr\}
\tag{A2}
\]

is assumed to be computable exactly.
\end{assumption}

\begin{assumption}[Stepsize]\label{ass:stepsize}
The stepsize $\eta$ satisfies  

\[
0<\eta\le\frac{1}{L}.
\tag{A3}
\]
\end{assumption}

\begin{assumption}[Existence of a minimiser]\label{ass:opt}
The optimal set $\mathcal{X}^\star\;=\;\{x\mid \Phi(x)=\inf_{z}\Phi(z)\}$ is non‑empty, and we denote $x^\star\in\mathcal{X}^\star$ an arbitrary optimal point.
\tag{A4}
\end{assumption}

--------------------------------------------------------------------

\section*{Main Convergence Theorem}
\begin{theorem}[Accelerated rate]\label{thm:rate}
Under Assumptions \ref{ass:smooth}–\ref{ass:opt}, the iterates $\{x_k\}$ generated by \eqref{2} satisfy for all $k\ge 1$

\[
\Phi(x_k)-\Phi(x^\star)\;\le\;
\frac{2\|x_0-x^\star\|^{2}}{\eta\,t_k^{2}}
\;\le\;
\frac{8L\|x_0-x^\star\|^{2}}{(k+1)^{2}} .
\tag{3}
\]

Consequently $\Phi(x_k)\to\Phi(x^\star)$ with the optimal $O(1/k^{2})$ rate for first‑order methods on the class of convex $L$‑smooth $f$ and convex $g$.
\end{theorem}

--------------------------------------------------------------------

\section*{Theoretical Properties}
\begin{enumerate}[label=\arabic*)]
\item \textbf{Stability.} The potential function  

\[
\mathcal{E}_k\;=\;t_k^{2}\bigl(\Phi(x_k)-\Phi(x^\star)\bigr)+\frac{1}{2\eta}\|x_k-x^\star\|^{2}
\]

is non‑increasing (see Lemma~\ref{lem:potential}), guaranteeing that the iterates remain bounded.

\item \textbf{Hyper‑parameter sensitivity.} The stepsize must respect $\eta\le1/L$; larger $\eta$ may violate the descent Lemma~\ref{lem:descent} and break the $O(1/k^2)$ guarantee. The momentum coefficients $\beta_k$ are uniquely determined by \eqref{1} and need not be tuned.

\item \textbf{Comparison to baselines.}
\begin{itemize}
\item \emph{Proximal Gradient (PG).} PG obtains $\Phi(x_k)-\Phi(x^\star)\le \dfrac{L\|x_0-x^\star\|^2}{2k}$.
\item \emph{APG-M} improves the dependence on $k$ from $1/k$ to $1/k^{2}$ while preserving the same per‑iteration complexity (one gradient and one proximal evaluation).
\end{itemize}
\end{enumerate}

--------------------------------------------------------------------

\section*{Discussion}
The theorem shows that adding Nesterov‑type momentum to the proximal gradient scheme yields the optimal accelerated rate for composite convex optimisation. The proof hinges on a careful combination of the smoothness descent lemma, the proximal inequality, and the algebraic structure of the sequence $\{t_k\}$.

--------------------------------------------------------------------

\section*{Preliminaries (Definitions and Lemmas)}
\begin{lemma}[Descent Lemma for $L$‑smooth $f$]\label{lem:descent}
For any $x,y\in\R^{d}$,
\[
f(y)\;\le\;f(x)+\langle\nabla f(x),y-x\rangle+\frac{L}{2}\|y-x\|^{2}.
\tag{4}
\]
\end{lemma}
\begin{proof}
Standard consequence of the $L$‑Lipschitz gradient property; see e.g. \cite{Nesterov2004}.
\end{proof}

\begin{lemma}[Proximal Inequality]\label{lem:prox}
For any $z\in\R^{d}$ and $u=\prox_{\eta g}(z)$,
\[
g(u)+\frac{1}{2\eta}\|u-z\|^{2}\;\le\;g(x)+\frac{1}{2\eta}\|x-z\|^{2}
\qquad\forall x\in\R^{d}.
\tag{5}
\]
\end{lemma}
\begin{proof}
Directly from the definition of the proximal operator as the minimiser of the right‑hand side.
\end{proof}

\begin{lemma}[Potential Decrease]\label{lem:potential}
Define $\mathcal{E}_k$ as in the discussion. Then for all $k\ge0$,
\[
\mathcal{E}_{k+1}\;\le\;\mathcal{E}_{k}.
\tag{6}
\]
\end{lemma}
\begin{proof}
The proof is a chain of algebraic manipulations; we present it in the main proof with explicit step numbers.
\end{proof}

--------------------------------------------------------------------

\section*{Main Proof (Step‑by‑Step Derivation)}
We prove Theorem~\ref{thm:rate}. Throughout the proof we fix an arbitrary optimal point $x^\star\in\mathcal{X}^\star$.

\noindent\textbf{Step 1 (Notation).} Set  

\[
s_k \;=\; y_k-\eta\nabla f(y_k),\qquad 
x_{k+1}= \prox_{\eta g}(s_k).
\tag{7}
\]

\noindent\textbf{Step 2 (Apply Lemma~\ref{lem:prox}).} For any $x$,
\[
g(x_{k+1})+\frac{1}{2\eta}\|x_{k+1}-s_k\|^{2}
\;\le\;
g(x)+\frac{1}{2\eta}\|x-s_k\|^{2}.
\tag{8}
\]

\noindent\textbf{Step 3 (Add $f$ and use Lemma~\ref{lem:descent}).}  
From \eqref{4} with $x=y_k$ and $y=x_{k+1}$,
\[
f(x_{k+1}) \le f(y_k)+\langle\nabla f(y_k),x_{k+1}-y_k\rangle
          +\frac{L}{2}\|x_{k+1}-y_k\|^{2}.
\tag{9}
\]

\noindent\textbf{Step 4 (Combine \eqref{8} and \eqref{9}).} Adding the two inequalities yields  

\[
\Phi(x_{k+1})\le
g(x)+f(y_k)+\langle\nabla f(y_k),x_{k+1}-y_k\rangle
+\frac{L}{2}\|x_{k+1}-y_k\|^{2}
+\frac{1}{2\eta}\|x-s_k\|^{2}
-\frac{1}{2\eta}\|x_{k+1}-s_k\|^{2}.
\tag{10}
\]

\noindent\textbf{Step 5 (Expand $s_k$).} Recall $s_k=y_k-\eta\nabla f(y_k)$. Hence  

\[
x_{k+1}-s_k = x_{k+1}-y_k+\eta\nabla f(y_k),\qquad
x-s_k = x-y_k+\eta\nabla f(y_k).
\tag{11}
\]

\noindent\textbf{Step 6 (Plug \eqref{11} into \eqref{10} and simplify).}  
A straightforward algebraic expansion gives  

\[
\begin{aligned}
\Phi(x_{k+1})\le
&g(x)+f(y_k)
+\frac{1}{2\eta}\|x-y_k\|^{2}
-\frac{1}{2\eta}\|x_{k+1}-y_k\|^{2}\\
&\quad +\underbrace{\Bigl\langle\nabla f(y_k),x_{k+1}-y_k\Bigr\rangle
-\langle\nabla f(y_k),x_{k+1}-y_k\rangle}_{=0}
+\frac{L-\frac{1}{\eta}}{2}\|x_{k+1}-y_k\|^{2}.
\end{aligned}
\tag{12}
\]

Since $\eta\le1/L$ (Assumption~\ref{ass:stepsize}), the coefficient $L-\frac{1}{\eta}\le0$, whence the last term is non‑positive and can be dropped:

\[
\Phi(x_{k+1})\le g(x)+f(y_k)+\frac{1}{2\eta}\|x-y_k\|^{2}
-\frac{1}{2\eta}\|x_{k+1}-y_k\|^{2}.
\tag{13}
\]

\noindent\textbf{Step 7 (Insert the convexity of $f$).} Convexity of $f$ gives  

\[
f(y_k)\le f(x)+\langle\nabla f(x),y_k-x\rangle .
\tag{14}
\]

Using \eqref{14} in \eqref{13} and rearranging,

\[
\Phi(x_{k+1})\le \Phi(x)+\langle\nabla f(x),y_k-x\rangle
+\frac{1}{2\eta}\|x-y_k\|^{2}
-\frac{1}{2\eta}\|x_{k+1}-y_k\|^{2}.
\tag{15}
\]

\noindent\textbf{Step 8 (Choose $x=x^\star$).} Because $x^\star$ minimises $\Phi$, $0\in\nabla f(x^\star)+\partial g(x^\star)$. In particular $\langle\nabla f(x^\star),y_k-x^\star\rangle\ge g(x^\star)-g(y_k)$. However we only need the inequality  

\[
\langle\nabla f(x^\star),y_k-x^\star\rangle\le f(y_k)-f(x^\star).
\tag{16}
\]

Applying \eqref{16} to \eqref{15} with $x=x^\star$ yields  

\[
\Phi(x_{k+1})- \Phi(x^\star)
\le
\frac{1}{2\eta}\|x^\star-y_k\|^{2}
-\frac{1}{2\eta}\|x_{k+1}-y_k\|^{2}.
\tag{17}
\]

\noindent\textbf{Step 9 (Express $y_k$ via momentum).} By definition of $y_k$ in \eqref{2},

\[
y_k - x^\star = x_k - x^\star + \beta_k (x_k - x_{k-1}).
\tag{18}
\]

Hence  

\[
\|x^\star-y_k\|^{2}
= \|x_k-x^\star\|^{2}+2\beta_k\langle x_k-x^\star,\,x_k-x_{k-1}\rangle
+\beta_k^{2}\|x_k-x_{k-1}\|^{2}.
\tag{19}
\]

Similarly  

\[
\|x_{k+1}-y_k\|^{2}
= \|x_{k+1}-x_k\|^{2}+2\beta_k\langle x_{k+1}-x_k,\,x_k-x_{k-1}\rangle
+\beta_k^{2}\|x_k-x_{k-1}\|^{2}.
\tag{20}
\]

\noindent\textbf{Step 10 (Form the potential).} Multiply both sides of \eqref{17} by $t_{k+1}^{2}$ and add $\frac{1}{2\eta}\|x_{k+1}-x^\star\|^{2}$ to both sides. After substituting \eqref{19}–\eqref{20} and using the identity  

\[
t_{k+1}^{2}\beta_k = t_k^{2}-t_{k+1},
\tag{21}
\]

which follows directly from the definition \eqref{1}, we obtain after a lengthy but elementary algebraic manipulation

\[
\mathcal{E}_{k+1}\;\le\;\mathcal{E}_{k},
\tag{22}
\]

where  

\[
\mathcal{E}_{k}\;=\;t_{k}^{2}\bigl(\Phi(x_{k})-\Phi(x^\star)\bigr)
+\frac{1}{2\eta}\|x_{k}-x^\star\|^{2}.
\tag{23}
\]

This establishes Lemma~\ref{lem:potential}.

\noindent\textbf{Step 11 (Bounding the potential).} Since $\mathcal{E}_{k}$ is non‑increasing and $\mathcal{E}_{0}= \frac{1}{2\eta}\|x_{0}-x^\star\|^{2}$, we have  

\[
t_{k}^{2}\bigl(\Phi(x_{k})-\Phi(x^\star)\bigr)
\le \frac{1}{2\eta}\|x_{0}-x^\star\|^{2}.
\tag{24}
\]

Rearranging gives  

\[
\Phi(x_{k})-\Phi(x^\star)\le
\frac{\|x_{0}-x^\star\|^{2}}{2\eta\,t_{k}^{2}}.
\tag{25}
\]

\noindent\textbf{Step 12 (Relating $t_k$ to $k$).} The recursion \eqref{1} yields the closed‑form lower bound  

\[
t_{k}\;\ge\;\frac{k+1}{2},
\tag{26}
\]

which can be proved by induction (see e.g. \cite{Nesterov2013}). Substituting \eqref{26} into \eqref{25} yields  

\[
\Phi(x_{k})-\Phi(x^\star)\le
\frac{2\|x_{0}-x^\star\|^{2}}{\eta\,(k+1)^{2}}
\;\le\;
\frac{8L\|x_{0}-x^\star\|^{2}}{(k+1)^{2}},
\tag{27}
\]

where we used $\eta\le 1/L$ in the last inequality. This is exactly the bound \eqref{3} claimed in Theorem~\ref{thm:rate}. ∎

--------------------------------------------------------------------

\section*{Rate Derivation (With Constants)}
The constant in \eqref{27) is $C=2/\eta$. If we select the largest admissible stepsize $\eta=1/L$, the bound becomes  

\[
\Phi(x_{k})-\Phi(x^\star)\le \frac{2L\|x_{0}-x^\star\|^{2}}{(k+1)^{2}}.
\tag{28}
\]

Thus the algorithm attains the optimal $O(L/k^{2})$ rate for the composite convex class.

--------------------------------------------------------------------

\section*{Special Cases}
\begin{enumerate}[label=\alph*)]
\item \textbf{Pure smooth case ($g\equiv0$).} The proximal step reduces to the identity, and the method coincides with Nesterov’s accelerated gradient, reproducing the classic $O(L/k^{2})$ bound.
\item \textbf{Indicator function $g=\iota_{\mathcal{C}}$.} The prox becomes Euclidean projection onto a convex set $\mathcal{C}$. The algorithm then yields the accelerated projected gradient method with the same rate.
\item \textbf{Strongly convex $f+g$.} If $\Phi$ is $\mu$‑strongly convex, a simple modification of the momentum coefficients (e.g., using $\beta_k=(\sqrt{L}-\sqrt{\mu})/(\sqrt{L}+\sqrt{\mu})$) leads to a linear convergence rate $O\bigl((1-\sqrt{\mu/L})^{k}\bigr)$.
\end{enumerate}

--------------------------------------------------------------------

\section*{References}
\begin{thebibliography}{9}
\bibitem{Nesterov2004}
Y.~Nesterov,
\textit{Introductory Lectures on Convex Optimization: A Basic Course},
Kluwer Academic Publishers, 2004.

\bibitem{Nesterov2013}
Y.~Nesterov,
\textit{Gradient methods for minimizing composite functions},
Mathematical Programming, 140(1):125–161, 2013.
\end{thebibliography}

\end{document}